\section{The extremal contradiction}\label{sec:incompatibility}

\begin{contract}
\textbf{Assumes:} \cref{prop:reduction}, i.e.\ a graph $G$ with positive indegrees satisfying
$\abs U+\abs{F^2(U)}<2\abs{F(U)}$ for every nonempty proper $U$. \quad
\textbf{Produces:} a contradiction (\cref{prop:no-concavity}), hence \cref{thm:main}.
\textbf{Why:} this is where the subset inequality and the pruning lemma meet. The subset inequality
applied to the rows and columns of two families of pairs yields two \emph{better} families that may
conflict; the pruning lemma repairs the conflicts without losing the improvement.
\end{contract}

\subsection{The setting}\label{sec:setting}

For the rest of the section fix an oriented graph $G$ with vertex set $X$, all indegrees positive,
and
\begin{equation}\label{eq:hypothesis}
 \abs U+\abs{F^2(U)}<2\abs{F(U)}
 \qquad\text{for every }\varnothing\ne U\subsetneq X .
\end{equation}
The object we extremise over is a \emph{pair of families of ordered pairs}, exactly as in
\cref{sec:pruning}. The two coordinate maps of \eqref{eq:Phi} act on them:
\begin{equation}\label{eq:Phi-again}
 \Phi_1 P=\{(i,j):(p,j)\in P,\ p\to i\},\qquad
 \Phi_2 Q=\{(i,j):(i,s)\in Q,\ s\to j\},
\end{equation}
and we abbreviate $\Phi_1^2P=\Phi_1(\Phi_1P)$ and $\Phi_2^2Q=\Phi_2(\Phi_2Q)$. Two families are
\emph{compatible} if $\Phi_1P\cap\Phi_2Q=\varnothing$ (\cref{rem:compat}). \Cref{fig:grid} fixes the
picture, and $\Delta=\{(v,v):v\in X\}$ is the diagonal.

\begin{figure}[t]
\centering
\begin{tikzpicture}[
  cell/.style={draw=black!35,minimum width=1.0cm,minimum height=0.82cm,inner sep=0pt},
  widecell/.style={draw=black!45,minimum width=1.5cm,minimum height=0.7cm,align=center,
                   font=\scriptsize,inner sep=1pt},
  filled/.style={cell,fill=blue!22,draw=blue!45!black},
  wfilled/.style={widecell,fill=blue!22,draw=blue!45!black},
  wghost/.style={widecell,fill=blue!7,draw=blue!35,dashed},
  rfilled/.style={widecell,fill=green!22,draw=green!40!black},
  rghost/.style={widecell,fill=green!7,draw=green!35,dashed},
  ar/.style={-{Stealth[length=1.6mm]},draw=black!75},
  lab/.style={font=\scriptsize,text=black!75},
  pan/.style={font=\scriptsize\bfseries,text=black!55}
]
% ---------------- (a) a family and one of its columns
\begin{scope}
\foreach \i in {0,1,2}{\foreach \j in {0,1,2}{
  \node[cell] at (\j*1.0+0.5, \i*0.82+0.41) {};}}
\node[filled] at (0.5,2.05) {};
\node[filled] at (1.5,2.05) {};
\node[filled] at (1.5,0.41) {};
\node[filled] at (2.5,1.23) {};
\draw[blue!55!black,dashed,thick] (1.02,0.0) rectangle (1.98,2.46);
\node[lab,anchor=north,text=blue!45!black] at (1.5,-0.06) {$P_j$};
\node[lab,anchor=north,align=center] at (1.5,-0.65) {a column};
\node[lab,anchor=north] at (1.5,-2.15) {(a) a family $P$};
\node[lab,anchor=east,rotate=90] at (-0.08,1.23) {rows $p$};
\node[lab,anchor=west] at (3.08,0.41) {$j$};

\end{scope}
% ---------------- (b) Phi_1 inside a column
\begin{scope}[shift={(4.7,0)}]
\node[wghost] (u) at (0,2.32) {$(i,j)$};
\node[wfilled] (c) at (0,1.23) {$(p,j)$};
\node[wghost] (d) at (0,0.14) {$(i',j)$};
\draw[ar] (c.north) -- (u.south);
\draw[ar] (c.south) -- (d.north);
\node[lab,anchor=west] at (0.9,1.78) {$p\to i$};
\node[lab,anchor=west] at (0.9,0.69) {$p\to i'$};
\node[lab,anchor=north,align=center] at (0,-0.65) {$\Phi_1P$: the\\ same column};
\node[lab,anchor=north] at (0,-2.15) {(b) $\Phi_1$};
\end{scope}
% ---------------- (c) Phi_2 inside a row
\begin{scope}[shift={(9.4,0)}]
\node[rghost] (l) at (-2.3,1.23) {$(i,j)$};
\node[rfilled] (c2) at (0,1.23) {$(i,s)$};
\node[rghost] (r2) at (2.3,1.23) {$(i,j')$};
\draw[ar] (c2.west) -- (l.east);
\draw[ar] (c2.east) -- (r2.west);
\node[lab,anchor=south] at (-1.15,1.62) {$s\to j$};
\node[lab,anchor=south] at (1.15,1.62) {$s\to j'$};
\node[lab,anchor=north,align=center] at (0,-0.65) {$\Phi_2Q$: the\\ same row};
\node[lab,anchor=north] at (0,-2.15) {(c) $\Phi_2$};
\end{scope}
\end{tikzpicture}

\caption{Families of pairs as cells of the grid $X\times X$. \emph{(a)} A family $P$; its
 \emph{columns} are the sets $P_j=\{p:(p,j)\in P\}$, one for each second coordinate (and its rows are
 defined the same way). \emph{(b)} The map $\Phi_1$ replaces the first coordinate $p$ by an
 out-neighbor $i$: it moves a cell within its column. \emph{(c)} The map $\Phi_2$ replaces the second
 coordinate $s$ by an out-neighbor $j$: it moves a cell within its row. \emph{Why it matters:} the
 hypothesis \eqref{eq:hypothesis} is a statement about \emph{sets of vertices}, so it can only be
 applied to the columns of $P$ and the rows of $Q$; summing it over those sets is what produces
 \eqref{eq:fiber-sums}.}
\label{fig:grid}
\end{figure}

\begin{proposition}\label{prop:no-concavity}
There is no nonempty finite oriented graph $G$ in which every vertex has an in-neighbor and whose
one-step image operator $F(U)=\{y: x\to y\text{ for some }x\in U\}$ satisfies
\begin{equation}\label{eq:impossible}
 \abs U+\abs{F^2(U)}<2\abs{F(U)}
 \qquad\text{for every }\varnothing\ne U\subsetneq V(G).
\end{equation}
\end{proposition}

\inwords{No graph satisfies the strict subset inequality at every proper nonempty set once every
vertex has an in-neighbor.}

\begin{ideabox}[title={The design of the extremal objective}]
We look for a compatible pair $(P,Q)$ that is as large as possible, and we pay for the size not only
with members but also with the points that the pair leaves uncovered. The objective is
\[
 M+u,\qquad M=\abs P+\abs Q,\qquad u=\abs X^2-\abs{\Phi_1P}-\abs{\Phi_2Q},
\]
where $u$ counts the points covered by neither family. Three considerations fix this choice.
\begin{enumerate}[leftmargin=1.3em,itemsep=2pt,topsep=1pt]
\item Only compatible pairs are admissible, because only then is $u$ the size of the complement of a
  disjoint union, and only then does the pruning lemma apply to the four sets we are about to build.
\item The subset inequality, summed over columns and rows, will produce families
  $\mathcal L=(X\times X)\setminus\Phi_2^2Q$ and $\mathcal R=(X\times X)\setminus\Phi_1^2P$ with
  $\abs{\mathcal L}+\abs{\mathcal R}>M+2u$: the gain is twice the uncovered count. The factor $2$ is
  what makes the objective $M+u$ rather than $M$: we must be able to pay for the pruning out of the
  gain and still be strictly ahead.
\item The pruning lemma deletes at most $\abs H+q\le u+q$ members, where $q$ is the number of points
  of $Z$ left uncovered after pruning, and the survivors leave $u'$ points uncovered in total. The
  new value is at least
  \[
   (\abs{\mathcal L}+\abs{\mathcal R})-(u+u')+u'=\abs{\mathcal L}+\abs{\mathcal R}-u>M+u,
  \]
  so the objective $M+u$ is exactly the quantity that the pruning cost cannot destroy. Including $u$
  in the objective is what lets the cost term $u'$ cancel against the refunded $u'$.
\end{enumerate}
\end{ideabox}

\begin{proof}[Proof of \cref{prop:no-concavity}]
Suppose that such a $G$ exists and put $X=V(G)$.

\steplabel{Step 1: the diagonal is compatible with itself.}
If $(i,j)\in\Phi_1\Delta\cap\Phi_2\Delta$, then $(p,j)\in\Delta$ with $p\to i$ forces $p=j$ and
$j\to i$, while $(i,s)\in\Delta$ with $s\to j$ forces $s=i$ and $i\to j$. Both $i\to j$ and $j\to i$
cannot hold in an oriented graph. So $\Phi_1\Delta\cap\Phi_2\Delta=\varnothing$, and the pair
$P=Q=\Delta$ is compatible. Among all compatible pairs with $\Delta\subseteq P\cap Q$ choose one
maximising
\begin{equation}\label{eq:objective}
 M+u,\qquad M=\abs P+\abs Q,\qquad u=\abs X^2-\abs{\Phi_1P}-\abs{\Phi_2Q};
\end{equation}
such a pair exists because $X$ is finite. Since the pair is compatible, $u$ is the number of points
lying in neither image.

\steplabel{Step 2: the columns and rows are nonempty and proper.}
For $j\in X$ let $P_j=\{p:(p,j)\in P\}$. It is nonempty, because $(j,j)\in\Delta\subseteq P$ gives
$j\in P_j$. It is also proper: if $P_j=X$, then $X\times\{j\}\subseteq\Phi_1P$, and choosing an
in-neighbor $w$ of $j$ (available by hypothesis) gives $(w,j)\in\Phi_2\Delta\subseteq\Phi_2Q$, which
contradicts compatibility. The same argument applies to the rows $Q_i=\{s:(i,s)\in Q\}$: $Q_i$ is
nonempty because $(i,i)\in Q$, and if $Q_i=X$ then $\{i\}\times X\subseteq\Phi_2Q$, while choosing
$w\to i$ gives $(i,w)\in\Phi_1\Delta\subseteq\Phi_1P$, again a contradiction.

\steplabel{Step 3: sum the subset inequality over the columns and rows.}
For each $j\in X$ the set $P_j$ is nonempty and proper, so \eqref{eq:hypothesis} gives
$\abs{P_j}+\abs{F^2(P_j)}<2\abs{F(P_j)}$. Summing over $j$ and using the fibrewise descriptions
\begin{equation}\label{eq:fibers}
 \abs P=\sum_{j}\abs{P_j},\qquad
 \Phi_1P\cap(X\times\{j\})=F(P_j)\times\{j\},\qquad
 \Phi_1^2P\cap(X\times\{j\})=F^2(P_j)\times\{j\},
\end{equation}
which hold because $\Phi_1$ acts inside each column, we obtain
\begin{equation}\label{eq:fiber-sums}
 \abs P+\abs{\Phi_1^2P}<2\abs{\Phi_1P}.
\end{equation}
The same argument with rows in place of columns gives
\begin{equation}\label{eq:fiber-sums-Q}
 \abs Q+\abs{\Phi_2^2Q}<2\abs{\Phi_2Q}.
\end{equation}

\steplabel{Step 4: take complements and gain $2u$.}
Define
\begin{equation}\label{eq:L-R}
 \mathcal L=(X\times X)\setminus\Phi_2^2Q,\qquad
 \mathcal R=(X\times X)\setminus\Phi_1^2P .
\end{equation}
Then, by \eqref{eq:fiber-sums} and \eqref{eq:fiber-sums-Q},
\begin{align}
 \abs{\mathcal L}+\abs{\mathcal R}
 &=2\abs X^2-\abs{\Phi_2^2Q}-\abs{\Phi_1^2P}\notag\\
 &>2\abs X^2+\bigl(\abs P+\abs Q\bigr)-2\bigl(\abs{\Phi_1P}+\abs{\Phi_2Q}\bigr)
 =M+2u ,
 \label{eq:gain}
\end{align}
where the last equality is \eqref{eq:objective}. So the two new families together have strictly more
than $M+2u$ members. \emph{The gain is twice the number of uncovered points}: the two strict
inequalities each contribute one copy of $u$, because each of them says that a first image is
strictly larger than the average of its family and its second image.

\steplabel{Step 5: the diagonal survives, and diagonals never conflict.}
\emph{$\Delta\subseteq\mathcal L\cap\mathcal R$.} Suppose $(v,v)\in\Phi_2^2Q$. Then there is $s$ with
$(v,s)\in\Phi_2Q$ and $s\to v$. But $(s,s)\in P$ and $s\to v$ give $(v,s)\in\Phi_1P$, so $(v,s)$
lies in both images, contradicting compatibility. Hence $(v,v)\notin\Phi_2^2Q$, i.e.\
$(v,v)\in\mathcal L$. Symmetrically, if $(v,v)\in\Phi_1^2P$ there is $p\to v$ with
$(p,v)\in\Phi_1P$; then $(p,p)\in Q$ and $p\to v$ give $(p,v)\in\Phi_2Q$, again a contradiction, so
$(v,v)\in\mathcal R$. \Cref{fig:diagonal} draws both arguments.

\begin{figure}[t]
\centering
\begin{tikzpicture}[
  cell/.style={draw=black!45,minimum width=1.6cm,minimum height=0.7cm,align=center,
               font=\scriptsize,inner sep=1pt},
  pc/.style={cell,fill=blue!14},
  qc/.style={cell,fill=green!14},
  bad/.style={cell,fill=red!20,draw=red!55!black},
  neut/.style={cell,fill=black!4},
  ar/.style={-{Stealth[length=1.6mm]},draw=black!70},
  lab/.style={font=\scriptsize,text=black!75}
]
% ---- case 1: (v,v) not in Phi_2^2 Q
\node[pc]   (ss) at (0,1.75) {$(s,s)\in P$};
\node[bad]  (vs) at (0,0)    {$(v,s)$};
\node[neut] (vv) at (2.6,0)  {$(v,v)$};
\draw[ar] (ss) -- node[right=2pt,lab] {$s\to v$} (vs);
\draw[ar] (vs) -- node[below=1pt,lab] {$s\to v$} (vv);
\node[font=\large,text=red!60!black,anchor=east] at (-1.0,0) {$\times$};
\node[lab,anchor=north,align=center] at (1.1,-0.6) {$(v,s)$ lies in $\Phi_1P$ and in $\Phi_2Q$};
\node[lab,anchor=south,align=center,text=black!55] at (2.6,1.75) {would be in $\Phi_2^2Q$};
% ---- case 2: (v,v) not in Phi_1^2 P
\node[qc]   (pp) at (6.4,0)     {$(p,p)\in Q$};
\node[bad]  (pv) at (9.0,0)     {$(p,v)$};
\node[neut] (vv2) at (9.0,-1.9) {$(v,v)$};
\draw[ar] (pp) -- node[above=1pt,lab] {$p\to v$} (pv);
\draw[ar] (pv) -- node[right=2pt,lab] {$p\to v$} (vv2);
\node[font=\large,text=red!60!black,anchor=south] at (9.0,0.55) {$\times$};
\node[lab,anchor=north,align=center] at (7.4,-2.7) {$(p,v)$ lies in $\Phi_1P$ and in $\Phi_2Q$};
\node[lab,anchor=west,align=left,text=black!55] at (10.0,-1.9) {would be in\\ $\Phi_1^2P$};
\end{tikzpicture}

\caption{Why the diagonal survives (Step~5 of that proof) and why diagonals never conflict. \emph{Look
 at} the two halves: if a diagonal cell were in $\Phi_2^2Q$ (left) or in $\Phi_1^2P$ (right), the
 cell marked $\times$ would lie in both images at once. \emph{Why it matters:} the pair
 $(\mathcal L,\mathcal R)$ may conflict, but the diagonal members are conflict-free, so the pruning
 lemma is allowed to keep them --- and the diagonal is what makes the pruned pair admissible, i.e.\
 a competitor to $(P,Q)$.}
\label{fig:diagonal}
\end{figure}

\emph{No diagonal member conflicts.} Let $(v,v)\in\Delta$ and $(i,s)\in\mathcal R$ conflict, i.e.\
$v\to i$ and $s\to v$. Then $(s,s)\in P$, and advancing its first coordinate twice along
$s\to v\to i$ puts $(i,s)$ in $\Phi_1^2P$, contradicting $(i,s)\in\mathcal R$. Symmetrically, if
$(p,j)\in\mathcal L$ conflicts with the diagonal member $(v,v)$, then $p\to v$ and $v\to j$, and
advancing the second coordinate of $(p,p)\in Q$ twice along $v\to j$ puts $(p,j)$ in $\Phi_2^2Q$,
contradicting $(p,j)\in\mathcal L$. So every diagonal member is conflict-free, and the preservation
clause of \cref{lem:pruning} applies to all of them.

\steplabel{Step 6: the collisions are uncovered.}
Let $Z,H$ be the sets attached to $\mathcal L,\mathcal R$ by \cref{def:conflict}: a conflict
\[
 (p,j)\in\mathcal L,\qquad (i,s)\in\mathcal R,\qquad p\to i,\qquad s\to j
\]
contributes $(p,s)$ to $H$ and $(i,j)$ to $Z$. We claim
\begin{equation}\label{eq:H-uncovered}
 H\subseteq X^2\setminus\bigl(\Phi_1P\cup\Phi_2Q\bigr),
 \qquad\text{hence}\qquad \abs H\le u .
\end{equation}
Fix such a conflict. If $(p,s)\in\Phi_1P$, then $(u,s)\in P$ for some $u\to p$, and the path
$u\to p\to i$ puts $(i,s)$ in $\Phi_1^2P$, contradicting $(i,s)\in\mathcal R$. If
$(p,s)\in\Phi_2Q$, then $(p,t)\in Q$ for some $t\to s$, and the path $t\to s\to j$ puts $(p,j)$ in
$\Phi_2^2Q$, contradicting $(p,j)\in\mathcal L$. So $(p,s)$ is in neither image, which is
\eqref{eq:H-uncovered}.

\steplabel{Step 7: prune, and contradict maximality.}
Apply \cref{lem:pruning} to $\mathcal L$ and $\mathcal R$. This gives conflict-free
$\mathcal L'\subseteq\mathcal L$ and $\mathcal R'\subseteq\mathcal R$ containing every initially
conflict-free member, so $\Delta\subseteq\mathcal L'\cap\mathcal R'$ by Step~5, and
\begin{equation}\label{eq:prune-cost}
 \abs{\mathcal L\setminus\mathcal L'}+\abs{\mathcal R\setminus\mathcal R'}\le\abs H+q
 \qquad\text{with}\qquad
 q=\abs{\{z\in Z:\ z\text{ covered by no member of }\mathcal L'\cup\mathcal R'\}} .
\end{equation}
Put
\begin{equation}\label{eq:u-prime}
 u'=\abs X^2-\abs{\Phi_1\mathcal L'}-\abs{\Phi_2\mathcal R'} .
\end{equation}
Since $\mathcal L'$ and $\mathcal R'$ have no conflicts, $\Phi_1\mathcal L'\cap\Phi_2\mathcal R'=\varnothing$
(\cref{rem:compat}), so $u'$ is the number of points covered by neither of them. A point is covered
by a kept left member exactly when it lies in $\Phi_1\mathcal L'$, and by a kept right member exactly
when it lies in $\Phi_2\mathcal R'$. Hence the points of $Z$ that are covered by no member of
$\mathcal L'\cup\mathcal R'$ are exactly the points of $Z$ outside $\Phi_1\mathcal L'\cup\Phi_2\mathcal R'$,
of which there are at most $u'$; so $q\le u'$. By \eqref{eq:prune-cost} and \eqref{eq:H-uncovered},
at most $\abs H+q\le u+u'$ members were deleted. Using \eqref{eq:gain},
\begin{align}
 \abs{\mathcal L'}+\abs{\mathcal R'}+u'
 &\ge\abs{\mathcal L}+\abs{\mathcal R}-(u+u')+u'\notag\\
 &=\abs{\mathcal L}+\abs{\mathcal R}-u>M+2u-u=M+u .
 \label{eq:contradiction}
\end{align}
But $(\mathcal L',\mathcal R')$ is an admissible pair: it is compatible and contains the diagonal. So
\eqref{eq:contradiction} contradicts the maximality of $(P,Q)$ in \eqref{eq:objective}. This proves
\cref{prop:no-concavity}.
\end{proof}

\begin{figure}[t]
\centering
\begin{tikzpicture}[
  stage/.style={draw=black!55,rounded corners=3pt,align=left,font=\scriptsize,
                inner sep=4pt,text width=10.4cm,fill=#1},
  stage/.default=white,
  val/.style={font=\scriptsize\bfseries,text=black!65,anchor=west},
  ar/.style={-{Stealth[length=2.2mm]},draw=black!70,thick},
  lab/.style={font=\scriptsize,text=black!70,align=left}
]
\node[stage=black!4] (s1) at (0,0)
 {\textbf{1. Extremal pair.} A compatible pair $(P,Q)$ with $\Delta\subseteq P\cap Q$ maximising
  $M+u$, where $M=\abs P+\abs Q$ and $u$ counts the points covered by neither family.};
\node[val] at (s1.east) {$M+u$};
\node[stage=blue!4,below=15mm of s1] (s2)
 {\textbf{2. Subset inequality.} Summed over the columns of $P$ and the rows of $Q$: for
  $\mathcal L=X^2\setminus\Phi_2^2Q$ and $\mathcal R=X^2\setminus\Phi_1^2P$,};
\node[val] at (s2.east) {$>M+2u$};
\node[stage=orange!5,below=15mm of s2] (s3)
 {\textbf{3. Pruning.} $\mathcal L,\mathcal R$ may conflict, but the diagonal members do not. Delete
  at most $\abs H+q\le u+u'$ members; the survivors leave $u'$ uncovered.};
\node[val] at (s3.east) {cost $\le u+u'$};
\node[stage=green!5,below=15mm of s3] (s4)
 {\textbf{4. Contradiction.} The survivors are compatible, contain $\Delta$, and are worth
  $>M+2u-u=M+u$.};
\node[val] at (s4.east) {impossible};
\draw[ar] (s1.south) -- node[left=3pt,lab] {$\Phi_1,\Phi_2$} (s2.north);
\draw[ar] (s2.south) -- node[left=3pt,lab] {prune (\cref{lem:pruning})} (s3.north);
\draw[ar] (s3.south) -- node[left=3pt,lab] {refund $u'$} (s4.north);
\end{tikzpicture}

\caption{The four moves of \cref{sec:incompatibility}, with the numbers attached. \emph{Read from
 left to right:} the extremal compatible pair $(P,Q)$ has value $M+u$; the subset inequalities,
 summed over columns and rows, produce the complements with value greater than $M+2u$; the pruning
 lemma deletes at most $u+u'$ members and leaves $u'$ points uncovered; the survivors are compatible,
 contain the diagonal, and have value greater than $M+u$ --- which is impossible. \emph{Why it
 matters:} the figure shows where each of the two copies of $u$ in the gain is spent: one pays for
 the collisions $H$, the other survives as the refunded $u'$.}
\label{fig:extremal}
\end{figure}

\subsection{Hypotheses, weights, and where strictness is needed}\label{sec:remarks4}

\begin{remark}[Where each hypothesis of \cref{prop:no-concavity} is used]\label{rem:hypotheses}
The proof uses exactly three properties of $G$, and it is worth knowing where.
\begin{itemize}[leftmargin=1.4em,itemsep=2pt,topsep=1pt]
\item \emph{The relation is oriented}, in four places: the diagonal is compatible with itself
 (Step~1); the diagonal survives both complement families (Step~5, first half); no diagonal member
 conflicts (Step~5, second half); and $N_2^+(v)=F^2(\{v\})\setminus F(\{v\})$ in the first place.
 All four uses are of the same fact --- there is no pair of opposite arcs --- and each of them would
 fail for a general digraph.
\item \emph{Every vertex has an in-neighbor}, in two places: a column $P_j$ that equals all of $X$
 would put $X\times\{j\}$ inside $\Phi_1P$ and a single in-neighbor of $j$ would then put $(w,j)$
 inside $\Phi_2Q$ (Step~2); and the same for rows. Without in-neighbors, \eqref{eq:hypothesis} cannot
 be applied to the columns and rows at all.
\item \emph{Strict concavity} \eqref{eq:hypothesis}, in one place only: Step~3, summed over the
 columns and the rows. Everything after that is counting.
\end{itemize}
So a reader checking the proof should watch exactly four steps: 1, 2, 3 and 5.
\end{remark}

\begin{remark}[Why the objective weights uncovered points with coefficient exactly $1$]\label{rem:lambda}
Suppose we maximised $\abs P+\abs Q+\lambda u$ for some other $\lambda>0$. The gain of Step~4 is
unchanged, $>M+2u$; the pruning deletes at most $\abs H+q\le u+u'$ members and the survivors have
$u'$ uncovered points, so the new value is larger than
\[
 M+2u-(u+u')+\lambda u'=\underbrace{M+\lambda u}_{\text{old value}}+(1-\lambda)u+(\lambda-1)u'
 =(\text{old value})+(\lambda-1)(u'-u).
\]
This is a strict improvement only if $(\lambda-1)(u'-u)\ge0$. For $\lambda=1$ the last term vanishes
and the strictness comes entirely from the strict gain, which is what we want; for $\lambda<1$ the
argument fails whenever $u'>u$, and for $\lambda>1$ it fails whenever $u'<u$. Since both cases occur,
$\lambda=1$ is the only weight that works for all instances. The coefficient $1$ on $u$ is therefore
forced by the structure of the pruning bound, not chosen for convenience.
\end{remark}

\begin{example}[Where the argument stops on a directed triangle]\label{ex:triangle}
Let $G$ be the directed triangle $a\to b\to c\to a$, which has all indegrees positive, and let
$P=Q=\Delta$. Then
\[
 \abs X^2=9,\qquad M=\abs P+\abs Q=6,\qquad u=9-3-3=3,
\]
because $\Phi_1\Delta=\Phi_2\Delta=\{(i,j):j\to i\}$, the arc set, has three elements, and
and $\Phi_1^2\Delta=\Phi_2^2\Delta=\Delta$, so the two complement families of Step~4 satisfy
\[
 \abs{\mathcal L}+\abs{\mathcal R}=2\cdot9-3-3=12=M+2u .
\]
The gain is $M+2u$ exactly, not more, because strict concavity fails for this graph: at $U=\{a\}$ we
have $1+1=2=2\cdot1$. Consequently the survivors are worth only
$\ge M+2u-(u+u')+u'=M+u$, which does not contradict maximality. \Cref{ex:triangle} therefore shows
precisely which strict inequality the whole argument rests on: the concavity hypothesis must be
strict, and the contradiction is the difference between `>' and `$\ge$'.
\end{example}

\begin{proof}[Proof of \cref{thm:main}]
If \cref{thm:main} were false, \cref{prop:reduction} would produce a nonempty finite oriented graph
with positive indegrees satisfying \eqref{eq:impossible} for every nonempty proper vertex set, and
\cref{prop:no-concavity} says that no such graph exists.
\end{proof}

\begin{checkpoint}
The proof of the main theorem is complete, and it used exactly two ingredients beyond the definitions:
the subset inequality (\cref{lem:deficit,lem:blowup}, giving \cref{prop:reduction}) and the pruning
lemma (\cref{lem:pruning}). Everything in the remaining two sections is a consequence.
\end{checkpoint}

\begin{remark}[Why the argument cannot be short-circuited]\label{rem:why-hard}
It is instructive to see which step fails if one tries to simplify the argument.
\begin{itemize}[leftmargin=1.4em,itemsep=2pt,topsep=1pt]
\item \emph{Why not just take $P=Q=\Delta$ and stop?} The diagonal is admissible --- that is Step~1
  --- but its value is $2\abs X+\abs X^2-2\abs A$ (because $\Phi_1\Delta=\{(i,j):j\to i\}$ and
  $\Phi_2\Delta=\{(i,j):i\to j\}$ are both copies of the arc set $A$), and what is needed is a pair
  that beats \emph{every} admissible pair. The diagonal's role is to make the maximum exist and to
  survive the pruning; the size has to come from the extremal choice.
\item \emph{Why not take larger $P,Q$ greedily?} Compatibility is a global condition: a single added
  member can cover an entire line of cells and destroy compatibility with the other family. The
  extremal choice is what handles this; there is no local augmentation argument.
\item \emph{Why is the pruning step not circular?} The pruning lemma is a statement about arbitrary
  relations; it does not know that $\mathcal L,\mathcal R$ come from an extremal pair. It is applied
  once, and the only thing it needs is that the diagonal members are conflict-free, which is Step~5.
\item \emph{Where does the strictness come from?} From \eqref{eq:hypothesis}, which is strict; the
  gain \eqref{eq:gain} is therefore strict, and the deletion bound is not strict, so the inequality
  survives. If \eqref{eq:hypothesis} were weak, the proof would only give $M+2u\ge\ldots$ and no
  contradiction.
\end{itemize}
\end{remark}
